\chapter{Station A: Spurensuche}\label{ch:spurensuche}

\begin{myremark}
	Der Unbekannte wusste, wo Lena wohnt. Dabei hat sie ihre Adresse nie gepostet.
	Wie ist das möglich? An dieser Station finden Sie heraus, welche Spuren Fotos, Apps und Geräte hinterlassen.
\end{myremark}

\section{Was in einem Foto steckt}

Jedes Foto, das Sie mit dem Handy aufnehmen, enthält neben dem Bild selbst zusätzliche Informationen, sogenannte \tib{Metadaten}.
Sie werden im \acrshort{EXIF}-Format direkt in der Bilddatei gespeichert, zum Beispiel:
\begin{itemize}
	\item Datum und Uhrzeit der Aufnahme,
	\item Hersteller und Modell des Handys,
	\item Einstellungen der Kamera (Blitz, Belichtung, \dots),
	\item und -- falls die Standortfunktion eingeschaltet ist -- die genauen \acrshort{GPS}-Koordinaten des Aufnahmeorts.
\end{itemize}

Viele grosse Plattformen entfernen die Standortdaten beim Hochladen. Das gilt aber nicht für alle Dienste, und oft nicht, wenn ein Bild als Datei verschickt wird (z.\,B. per E-Mail, AirDrop oder \enquote{als Dokument}).
Ausserdem verrät häufig schon der Bildinhalt viel: ein Strassenschild, eine Hausnummer, das Schullogo auf dem Pullover.

\begin{myexercise}[Mein eigenes Foto]
	Öffnen Sie auf Ihrem Handy ein Foto, das Sie selbst aufgenommen haben, und lassen Sie sich die Informationen dazu anzeigen
	(iPhone: Foto antippen und nach oben wischen; Android/Google Fotos: Foto öffnen, Menü \enquote{\dots} oder nach oben wischen).
	Notieren Sie, was Sie finden. Zeigen Sie das Foto niemandem.

	\begin{center}
		\renewcommand{\arraystretch}{1.6}
		\begin{tabular}{p{5cm}p{8cm}}
			\textbf{Information}       & \textbf{Gefunden? Was steht dort?} \\
			\midrule
			Datum und Uhrzeit          &                                    \\
			Gerät (Modell)             &                                    \\
			Ort / Karte                &                                    \\
			Weitere Angaben            &                                    \\
		\end{tabular}
	\end{center}
\end{myexercise}

\begin{myexercise}[Detektivarbeit am Testfoto]\label{ex:testfoto}
	Laden Sie das Testfoto \enquote{Lena\_Garten.jpg} herunter, das Ihre Lehrperson bereitgestellt hat.
	Öffnen Sie die Dateiinformationen auf dem Laptop
	(Windows: Rechtsklick, \enquote{Eigenschaften}, Register \enquote{Details}; macOS: in \enquote{Vorschau} öffnen, \enquote{Werkzeuge}, \enquote{Informationen einblenden}).
	\begin{enumerate}
		\item Wann wurde das Foto aufgenommen?
		\item Mit welchem Gerät?
		\item Wo wurde es aufgenommen? Geben Sie die Koordinaten in eine Kartenanwendung ein.
		\item Erklären Sie in einem Satz, wie der Unbekannte Lenas Wohnort herausfinden konnte.
	\end{enumerate}
	\begin{myanswer}
		(a) 30.~Januar 2027, 14:23~Uhr. (b) Apple iPhone~13.
		(c) $47{,}35455^\circ$\,N, $8{,}55127^\circ$\,E: Chinagarten in Zürich (im Fall steht dieser Ort für Lenas Garten).
		(d) Lena hat ein Foto aus ihrem Garten verschickt, das die \acrshort{GPS}-Koordinaten enthielt. Wer die Datei erhielt, konnte den Ort auf einer Karte anzeigen.
	\end{myanswer}
\end{myexercise}

\begin{myexercise}[Standort beim Fotografieren ausschalten]
	Suchen Sie auf Ihrem Handy die Einstellung, mit der die Kamera \emph{keinen} Standort mehr speichert.
	Notieren Sie den Weg dorthin (z.\,B. \enquote{Einstellungen $\rightarrow$ \dots}).
	Entscheiden Sie selbst, ob Sie die Einstellung ändern möchten.
	\schreibzeile
	\schreibzeile
	\begin{myanswer}
		iPhone: Einstellungen $\rightarrow$ Datenschutz \& Sicherheit $\rightarrow$ Ortungsdienste $\rightarrow$ Kamera $\rightarrow$ \enquote{Nie}.
		Android: in der Kamera-App unter Einstellungen die Option \enquote{Standort speichern} bzw. \enquote{Standort-Tags} ausschalten (Bezeichnung je nach Hersteller).
		Beim Teilen auf dem iPhone kann der Ort zudem unter \enquote{Optionen} für ein einzelnes Bild entfernt werden.
	\end{myanswer}
\end{myexercise}

\section{Was Apps über Sie wissen dürfen}

Apps fragen um \tib{Berechtigungen}: Zugriff auf Standort, Kamera, Mikrofon, Kontakte oder Fotos.
Manche brauchen diese wirklich (eine Karten-App den Standort), andere sammeln damit vor allem Daten für Werbung.

\begin{myexercise}[Berechtigungs-Check]
	Wählen Sie drei Apps auf Ihrem Handy aus und prüfen Sie deren Berechtigungen
	(iPhone: Einstellungen $\rightarrow$ App auswählen; Android: Einstellungen $\rightarrow$ Apps $\rightarrow$ App $\rightarrow$ Berechtigungen).
	Tragen Sie ein, was die App darf (\faCheck) und ob sie das für ihre Aufgabe wirklich braucht.

	\begin{center}
		\renewcommand{\arraystretch}{1.6}
		\small
		\begin{tabular}{p{2.4cm}cccc p{3.2cm}}
			\textbf{App} & \textbf{Standort} & \textbf{Kamera} & \textbf{Mikrofon} & \textbf{Kontakte} & \textbf{Braucht sie das?} \\
			\midrule
			             &                   &                 &                   &                   &                           \\
			             &                   &                 &                   &                   &                           \\
			             &                   &                 &                   &                   &                           \\
		\end{tabular}
	\end{center}
\end{myexercise}

\section{Ihr Gerät hat einen Fingerabdruck}

Viele Menschen lehnen auf Webseiten die Cookies ab, damit sie nicht verfolgt werden.
Die Werbebranche nutzt deshalb zunehmend eine andere Methode: das \tib{Fingerprinting}.
Dabei liest eine Webseite Dutzende harmlose Merkmale Ihres Geräts aus -- Bildschirmgrösse, Zeitzone, Sprache, installierte Schriftarten, Version des Betriebssystems und vieles mehr.
Einzeln sagen diese Merkmale wenig aus. Zusammen ergeben sie aber eine Kombination, die je nach Messung bei 80 bis 90 Prozent aller Geräte einzigartig ist.
So kann ein Gerät wiedererkannt werden, auch ohne Cookies und ohne dass Sie es bemerken.

\begin{myexercise}[Wie einzigartig ist mein Browser?]
	Öffnen Sie auf Ihrem Laptop oder Handy die Seite \url{https://coveryourtracks.eff.org} und starten Sie den Test (\enquote{Test your browser}).
	\begin{enumerate}
		\item Hat Ihr Browser einen einzigartigen Fingerabdruck (\enquote{unique fingerprint})?
		\item Nennen Sie drei Merkmale, die der Test über Ihr Gerät herausgefunden hat.
	\end{enumerate}
	\schreibzeile
	\schreibzeile
\end{myexercise}

\begin{mychallenge}[Inkognito hilft nicht]
	Wiederholen Sie den Test in einem privaten Fenster (Inkognito-Modus).
	Erklären Sie, warum der Inkognito-Modus zwar den Verlauf und die Cookies löscht, aber kaum gegen Fingerprinting schützt.
	\schreibzeile
	\schreibzeile
	\begin{myanswer}
		Der Inkognito-Modus speichert lokal keinen Verlauf und löscht Cookies nach dem Schliessen.
		Die Merkmale des Geräts (Bildschirm, Schriftarten, Zeitzone, Betriebssystem, \dots) bleiben aber gleich und werden der Webseite weiterhin mitgeteilt.
		Darum ergibt sich derselbe oder ein sehr ähnlicher Fingerabdruck.
	\end{myanswer}
\end{mychallenge}

\begin{mychallenge}[Spurensuche im Bild]
	Nicht nur Metadaten verraten einen Ort. Suchen Sie in einem eigenen Foto, das Sie gerne posten würden, nach Hinweisen im Bildinhalt
	(Schilder, Hausnummern, Logos, Spiegelungen). Was müssten Sie abdecken oder wegschneiden?
\end{mychallenge}

\begin{fallakte}[-- Station A]
	Was ist bei Lena schiefgelaufen? Was hätte sie anders machen können?
	\schreibzeile
	\schreibzeile
	\schreibzeile
\end{fallakte}
